\documentclass[10pt,letterpaper]{article}
\usepackage[utf8]{inputenc}
\usepackage[intlimits]{amsmath}
\usepackage{amsfonts}
\usepackage{amssymb}
\usepackage{ragged2e}
\usepackage[letterpaper, margin=1in]{geometry}
\usepackage{graphicx}
\usepackage{cancel}
\usepackage{mathtools}
\usepackage{tabularx}
\usepackage{arydshln}
\usepackage{tensor}
\usepackage{array}
\usepackage{xcolor}
\usepackage[boxed]{algorithm}
\usepackage[noend]{algpseudocode}
\usepackage{listings}
\usepackage{textcomp}
% \usepackage[pdf,tmpdir,singlefile]{graphviz}
\usepackage{mathrsfs}
\usepackage{bbm}
\usepackage{tikz}
\usepackage{tikz-cd}
\usepackage{enumitem}
\usepackage{arydshln}
\usepackage{relsize}
\usepackage{multicol}
\usepackage{scalerel}

\usetikzlibrary{bayesnet}
\setlist{noitemsep}

%%%%%%%%%%%%%%%%%%%%%%%%%%%%%
% Formatting commands
%%%%%%%%%%%%%%%%%%%%%%%%%%%%%
\newcommand{\n}{\hfill\break}
\newcommand{\up}{\vspace{-\baselineskip}}
\newcommand{\hangblock}[2]{\par\noindent\settowidth{\hangindent}{\textbf{#1: }}\textbf{#1: }\!\!\!#2}
\newcommand{\lemma}[1]{\hangblock{Lemma}{#1}}
\newcommand{\defn}[1]{\hangblock{Defn}{#1}}
\newcommand{\thm}[1]{\hangblock{Thm}{#1}}
\newcommand{\cor}[1]{\hangblock{Cor}{#1}}
\newcommand{\prop}[1]{\hangblock{Prop}{#1}}
\newcommand{\ex}[1]{\hangblock{Ex}{#1}}
\newcommand{\exer}[1]{\hangblock{Exer}{#1}}
\newcommand{\fact}[1]{\hangblock{Fact}{#1}}
\newcommand{\remark}[1]{\hangblock{Remark}{#1}}
\newcommand{\proven}{\;$\square$\n}
\newcommand{\problem}[1]{\par\noindent{#1}\n}
\newcommand{\problempart}[2]{\par\noindent\indent{}\settowidth{\hangindent}{\textbf{(#1)} \indent{}}\textbf{(#1)} #2\n}
\newcommand{\ptxt}[1]{\textrm{\textnormal{#1}}}
\newcommand{\inlineeq}[1]{\centerline{$\displaystyle #1$}}
\newcommand{\pageline}{\noindent\rule{\textwidth}{0.1pt}}

%%%%%%%%%%%%%%%%%%%%%%%%%%%%%
% Math commands
%%%%%%%%%%%%%%%%%%%%%%%%%%%%%
% Set Theory
\newcommand{\card}[1]{\left|#1\right|}
\newcommand{\set}[1]{\left\{#1\right\}}
\newcommand{\setmid}{\;\middle|\;}
\newcommand{\ps}[1]{\mathcal{P}\left(#1\right)}
\newcommand{\pfinite}[1]{\mathcal{P}^{\ptxt{finite}}\left(#1\right)}
\newcommand{\naturals}{\mathbb{N}}
\newcommand{\N}{\naturals}
\newcommand{\integers}{\mathbb{Z}}
\newcommand{\Z}{\integers}
\newcommand{\rationals}{\mathbb{Q}}
\newcommand{\Q}{\rationals}
\newcommand{\reals}{\mathbb{R}}
\newcommand{\R}{\reals}
\newcommand{\complex}{\mathbb{C}}
\newcommand{\C}{\complex}
\newcommand{\halfPlane}{\mathbb{H}}
\let\H\relax
\newcommand{\H}{\halfPlane}
\newcommand{\comp}{^{\complement}}
\DeclareMathOperator{\Hom}{Hom}
\newcommand{\Ind}{\mathbbm{1}}
\newcommand{\cut}{\setminus}
\DeclareMathOperator{\elem}{elem}

% Graph Theory
\let\deg\relax
\DeclareMathOperator{\deg}{deg}
\newcommand{\degp}{\ptxt{deg}^{+}}
\newcommand{\degn}{\ptxt{deg}^{-}}
\newcommand{\precdot}{\mathrel{\ooalign{$\prec$\cr\hidewidth\hbox{$\cdot\mkern0.5mu$}\cr}}}
\newcommand{\succdot}{\mathrel{\ooalign{$\cdot\mkern0.5mu$\cr\hidewidth\hbox{$\succ$}\cr\phantom{$\succ$}}}}
\DeclareMathOperator{\cl}{cl}
\DeclareMathOperator{\affdim}{affdim}

% Probability
\newcommand{\parSymbol}{\P}
\newcommand{\Prob}{\mathbb{P}}
\renewcommand{\P}{\Prob}
\newcommand{\Avg}{\mathbb{E}}
\newcommand{\E}{\Avg}
\DeclareMathOperator{\Var}{Var}
\DeclareMathOperator{\cov}{cov}
\DeclareMathOperator{\Unif}{Unif}
\DeclareMathOperator{\Binom}{Binom}
\newcommand{\CI}{\mathrel{\text{\scalebox{1.07}{$\perp\mkern-10mu\perp$}}}}

% Standard Math
\newcommand{\inv}{^{-1}}
\newcommand{\abs}[1]{\left|#1\right|}
\newcommand{\ceil}[1]{\left\lceil{}#1\right\rceil{}}
\newcommand{\floor}[1]{\left\lfloor{}#1\right\rfloor{}}
\newcommand{\conj}[1]{\overline{#1}}
\newcommand{\of}{\circ}
\newcommand{\tri}{\triangle}
\newcommand{\inj}{\hookrightarrow}
\newcommand{\surj}{\twoheadrightarrow}
\newcommand{\ndiv}{\nmid}
\renewcommand{\epsilon}{\varepsilon}
\newcommand{\divides}{\mid}
\newcommand{\ndivides}{\nmid}
\DeclareMathOperator{\lcm}{lcm}
\DeclareMathOperator{\sgn}{sgn}
\newcommand{\map}[4]{\!\!\!\begin{array}[t]{rcl}#1 & \!\!\!\!\to & \!\!\!\!#2\\ #3 & \!\!\!\!\mapsto & \!\!\!\!#4\end{array}}
\newcommand{\bigsum}[2]{\smashoperator[lr]{\sum_{\scalebox{#1}{$#2$}}}}

% Linear Algebra
\newcommand{\Id}{\textrm{\textnormal{Id}}}
\newcommand{\im}{\textrm{\textnormal{im}}}
\newcommand{\norm}[1]{\abs{\abs{#1}}}
\newcommand{\tpose}{^{T}\!}
\newcommand{\iprod}[1]{\left<#1\right>}
\DeclareMathOperator{\tr}{tr}
\DeclareMathOperator{\trace}{tr}
\newcommand{\chgBasMat}[3]{\!\!\tensor*[_{#1}]{\left[#2\right]}{_{#3}}}
\newcommand{\vecBas}[2]{\tensor*[]{\left[#1\right]}{_{#2}}}
\DeclareMathOperator{\GL}{GL}
\DeclareMathOperator{\Mat}{Mat}
\DeclareMathOperator{\vspan}{span}
\DeclareMathOperator{\rank}{rank}
\newcommand{\V}[1]{\vec{#1}}
\DeclareMathOperator{\proj}{proj}
\DeclareMathOperator{\compProj}{comp}
\DeclareMathOperator{\row}{row}
\newcommand{\smallPMatrix}[1]{\paren{\begin{smallmatrix}#1\end{smallmatrix}}}
\newcommand{\smallBMatrix}[1]{\brack{\begin{smallmatrix}#1\end{smallmatrix}}}
\newcommand{\pmat}[1]{\begin{pmatrix}#1\end{pmatrix}}
\newcommand{\bmat}[1]{\begin{bmatrix}#1\end{bmatrix}}

% Multilinear Algebra
\newcommand{\Lsym}{\L}
\let\L\relax
\DeclareMathOperator{\L}{\mathscr{L}}
\DeclareMathOperator{\A}{\mathcal{A}}
\DeclareMathOperator{\Alt}{Alt}
\DeclareMathOperator{\Sym}{Sym}
\newcommand{\ot}{\otimes}
\newcommand{\ox}{\otimes}
\DeclareMathOperator{\asc}{asc}
\DeclareMathOperator{\asSet}{set}
\DeclareMathOperator{\sort}{sort}
\DeclareMathOperator{\ringA}{\mathring{A}}

% Topology
\newcommand{\closure}[1]{\overline{#1}}
\newcommand{\uball}{\mathcal{U}}
\DeclareMathOperator{\Int}{Int}
\DeclareMathOperator{\Ext}{Ext}
\DeclareMathOperator{\Bd}{Bd}
\DeclareMathOperator{\rInt}{rInt}
\DeclareMathOperator{\ch}{ch}
\DeclareMathOperator{\ah}{ah}
\newcommand{\LargerTau}{\mathlarger{\mathlarger{\mathlarger{\mathlarger{\tau}}}}}
\newcommand{\Tau}{\mathcal{T}}

% Analysis
\DeclareMathOperator{\Graph}{Graph}
\DeclareMathOperator{\epi}{epi}
\DeclareMathOperator{\hypo}{hypo}
\DeclareMathOperator{\supp}{supp}
\newcommand{\lint}[2]{\underset{#1}{\overset{#2}{{\color{black}\underline{{\color{white}\overline{{\color{black}\int}}\color{black}}}}}}}
\newcommand{\uint}[2]{\underset{#1}{\overset{#2}{{\color{white}\underline{{\color{black}\overline{{\color{black}\int}}\color{black}}}}}}}
\newcommand{\alignint}[2]{\underset{#1}{\overset{#2}{{\color{white}\underline{{\color{white}\overline{{\color{black}\int}}\color{black}}}}}}}
\newcommand{\extint}{\ptxt{ext}\int}
\newcommand{\extalignint}[2]{\ptxt{ext}\alignint{#1}{#2}}
\newcommand{\conv}{\ast}
\newcommand{\pd}[2]{\frac{\partial{}#1}{\partial{}#2}}
\newcommand{\del}{\nabla}
\DeclareMathOperator{\grad}{grad}
\DeclareMathOperator{\curl}{curl}
\let\div\relax
\DeclareMathOperator{\div}{div}
\DeclareMathOperator{\vol}{vol}

% Complex Analysis
\let\Re\relax
\DeclareMathOperator{\Re}{Re}
\let\Im\relax
\DeclareMathOperator{\Im}{Im}
\DeclareMathOperator{\Res}{Res}

% Abstract Algebra
\DeclareMathOperator{\ord}{ord}
\newcommand{\generated}[1]{\left<#1\right>}
\newcommand{\cycle}[1]{\smallPMatrix{#1}}
\newcommand{\id}{\ptxt{id}}
\newcommand{\iso}{\cong}
\DeclareMathOperator{\Aut}{Aut}
\DeclareMathOperator{\SL}{SL}
\DeclareMathOperator{\op}{op}
\newcommand{\isom}[4]{\!\!\!\begin{array}[t]{rcl}#1 & \!\!\!\!\overset{\sim}{\to} & \!\!\!\!#2\\ #3 & \!\!\!\!\mapsto & \!\!\!\!#4\end{array}}
\newcommand{\F}{\mathbb{F}}
\newcommand{\acts}{\;\rotatebox[origin=c]{-90}{$\circlearrowright$}\;}
\newcommand{\disjunion}{\mathrel{\text{\scalebox{1.07}{$\perp\mkern-10mu\perp$}}}}

% Convex Optimization
\newcommand{\sectionSymbol}{\S}
\let\S\relax
\newcommand{\S}{\mathbb{S}}
\DeclareMathOperator{\dist}{dist}
\DeclareMathOperator{\dom}{dom}
\DeclareMathOperator{\diag}{diag}
\DeclareMathOperator{\ones}{\mathbbm{1}}
\newcommand{\minimizeOver}[3]{\begin{array}{rl}\underset{#1}{\ptxt{minimize}} & #2\\ \ptxt{subject to} & #3\end{array}}
\newcommand{\maximizeOver}[3]{\begin{array}{rl}\underset{#1}{\ptxt{maximize}} & #2\\ \ptxt{subject to} & #3\end{array}}
\newcommand{\minimizationProblem}[2]{\minimizeOver{}{#1}{#2}}
\newcommand{\maximizationProblem}[2]{\maximizeOver{}{#1}{#2}}

% Proofs
\newcommand{\st}{s.t.}
\newcommand{\unique}{!}
\newcommand{\iffdef}{\overset{\ptxt{def}}{\Leftrightarrow}}
\newcommand{\eqdef}{\overset{\ptxt{def}}{=}}
\newcommand{\eqVertical}{\rotatebox[origin=c]{90}{=}}
\newcommand{\mapsfrom}{\mathrel{\reflectbox{\ensuremath{\mapsto}}}}
\newcommand{\mapsdown}{\rotatebox[origin=c]{-90}{$\mapsto$}\mkern2mu}
\newcommand{\mapsup}{\rotatebox[origin=c]{90}{$\mapsto$}\mkern2mu}
\newcommand{\from}{\!\mathrel{\reflectbox{\ensuremath{\to}}}}

% Brackets
\newcommand{\paren}[1]{\left(#1\right)}
\renewcommand{\brack}[1]{\left[#1\right]}
\renewcommand{\brace}[1]{\left\{#1\right\}}
\newcommand{\ang}[1]{\left<#1\right>}

% Algorithms
\algrenewcommand{\algorithmiccomment}[1]{\hskip 1em \texttt{// #1}}
\algrenewcommand\algorithmicrequire{\textbf{Input:}}
\algrenewcommand\algorithmicensure{\textbf{Output:}}
\newcommand{\algP}{\ptxt{\textbf{P}}}
\newcommand{\algNP}{\ptxt{\textbf{NP}}}
\newcommand{\algNPC}{\ptxt{\textbf{NP-Complete}}}
\newcommand{\algNPH}{\ptxt{\textbf{NP-Hard}}}
\newcommand{\algEXP}{\ptxt{\textbf{EXP}}}

%%%%%%%%%%%%%%%%%%%%%%%%%%%%%
% Other commands
%%%%%%%%%%%%%%%%%%%%%%%%%%%%%
\newcommand{\flag}[1]{\textbf{\textcolor{red}{#1}}}
\newcommand{\uSym}{\u}
\let\u\relax
\newcommand{\u}[1]{\underline{#1}}
\newcommand{\bSym}{\b}
\let\b\relax
\newcommand{\b}[1]{\textbf{#1}}
\newcommand{\iSym}{\i}
\let\i\relax
\newcommand{\i}[1]{\textit{#1}}

%%%%%%%%%%%%%%%%%%%%%%%%%%%%%%%%%%%%%%%
% Make l's curvy in math environments %
%%%%%%%%%%%%%%%%%%%%%%%%%%%%%%%%%%%%%%%
\mathcode`l="8000
\begingroup
\makeatletter
\lccode`\~=`\l
\DeclareMathSymbol{\lsb@l}{\mathalpha}{letters}{`l}
\lowercase{\gdef~{\ifnum\the\mathgroup=\m@ne \ell \else \lsb@l \fi}}%
\endgroup

%%%%%%%%%%%%%%%%%%%%%%%%%
% Fix \vdots and \ddots %
%%%%%%%%%%%%%%%%%%%%%%%%%
\usepackage{letltxmacro}
\LetLtxMacro\orgvdots\vdots
\LetLtxMacro\orgddots\ddots

\makeatletter
\DeclareRobustCommand\vdots{%
	\mathpalette\@vdots{}%
}
\newcommand*{\@vdots}[2]{%
	% #1: math style
	% #2: unused
	\sbox0{$#1\cdotp\cdotp\cdotp\m@th$}%
	\sbox2{$#1.\m@th$}%
	\vbox{%
		\dimen@=\wd0 %
		\advance\dimen@ -3\ht2 %
		\kern.5\dimen@
		% remove side bearings
		\dimen@=\wd2 %
		\advance\dimen@ -\ht2 %
		\dimen2=\wd0 %
		\advance\dimen2 -\dimen@
		\vbox to \dimen2{%
			\offinterlineskip
			\copy2 \vfill\copy2 \vfill\copy2 %
		}%
	}%
}
\DeclareRobustCommand\ddots{%
	\mathinner{%
		\mathpalette\@ddots{}%
		\mkern\thinmuskip
	}%
}
\newcommand*{\@ddots}[2]{%
	% #1: math style
	% #2: unused
	\sbox0{$#1\cdotp\cdotp\cdotp\m@th$}%
	\sbox2{$#1.\m@th$}%
	\vbox{%
		\dimen@=\wd0 %
		\advance\dimen@ -3\ht2 %
		\kern.5\dimen@
		% remove side bearings
		\dimen@=\wd2 %
		\advance\dimen@ -\ht2 %
		\dimen2=\wd0 %
		\advance\dimen2 -\dimen@
		\vbox to \dimen2{%
			\offinterlineskip
			\hbox{$#1\mathpunct{.}\m@th$}%
			\vfill
			\hbox{$#1\mathpunct{\kern\wd2}\mathpunct{.}\m@th$}%
			\vfill
			\hbox{$#1\mathpunct{\kern\wd2}\mathpunct{\kern\wd2}\mathpunct{.}\m@th$}%
		}%
	}%
}
\makeatother

\newcommand{\B}{
	\begin{tikzpicture}
	\filldraw [fill=red, draw=black] (0, 0) rectangle (0.37, 0.45);
	\draw [line width=0.5mm, white ] (0.1,0.08) -- (0.1,0.38);
	\draw[line width=0.5mm, white ] (0.1, 0.35) .. controls (0.2, 0.35) and (0.4, 0.2625) .. (0.1, 0.225);
	\draw[line width=0.5mm, white ] (0.1, 0.225) .. controls (0.2, 0.225) and (0.4, 0.1625) .. (0.1, 0.1);
	\end{tikzpicture}
}

\author{Thomas Cohn}
\title{Math 493 Lecture 12}
\date{10/16/19} % Can also use \today

\begin{document}
\maketitle
\setlength\RaggedRightParindent{\parindent}
\RaggedRight

\defn{
	Let $G$ be a group, and $X$ be a set. An \textbf{action} of $G$ on $X$ is a function $\map{G\times{}X}{X}{(g,x)}{g\cdot{}x}$.\n
	\st{}
	\begin{enumerate}[leftmargin=4\parindent]
		\item $1\cdot{}x=x$, $\forall{}x\in{}X$
		\item $(gh)x=g(hx)$, $\forall{}g,h\in{}G$, $\forall{}x\in{}X$
	\end{enumerate}\up\n
	We may also say $X$ is a \textbf{$G$-set}.\n
}

\ex{
	\begin{enumerate}
		\item $M$, the group of rigid motions of the plane $P=\R^{2}$. $M$ acts on $P$.
		\item $G$, any group, acts on $X=G$ by
		\begin{enumerate}
			\item $g\cdot{}h=gh$
			\item $g\cdot{}h=hg\inv$. Check:\n
			$\forall{}g,h,x\in{}G$, $(gh)\cdot{}x=x(gh)\inv=xh\inv{}g\inv=(xh\inv)g\inv=g\cdot(xh\inv)=g\cdot(h\cdot{}x)$
			\item $g\cdot{}x=gxg\inv$
		\end{enumerate}
		\item $G=S_{n}$, $X=\set{1,2,\ldots,n}$. $\sigma\in{}G$, $i\in{}X$, $\sigma\cdot{}i=\sigma(i)$.
		\item $G=\GL_{n}(F)$, ($F$ is a field). $X=F^{n}$. $g\cdot{}x$ is just matrix-vector multiplication.
	\end{enumerate}
}

\defn{
	$X$ and $Y$ are $G$-sets. A \textbf{homomorphism} from $X$ to $Y$ is a function $f:X\to{}Y$ such that $f(gx)=gf(x)$, $\forall{}x\in{}X,g\in{}G$. We say that $f$ is a \textbf{$G$-map}.\n
}

\defn{
	An \textbf{isomorphism} of $G$-sets is a bijective homomorphism.\n
}

\defn{
	If $G$ is any group, $X$ is any set, we have the \textbf{trivial action} of $G$ on $X$ by $g\cdot{}x=x$, $\forall{}g\in{}G,x\in{}X$.\n
}

\defn{
	Let $X$ be a $G$-set. The \textbf{orbit} of $x\in{}X$ (under $G$) is $O_{x}=\set{g\cdot{}x\mid{}g\in{}G}$.\n
}

\par\noindent
Note: we can define an equivalence relation on $X$ by $x\sim{}y$ if $\exists{}g\in{}G$ \st{} $x=gy$. $O_{x}$ is the equivalence class.\n
Equivalence classes partition their set, so the orbits partition $X$.\n

\par\noindent
Observation: given $y\in{}O_{x}g\in{}G$, $gy\in{}O_{x}$, because $y=hx$ for some $h\in{}O_{x}$, so $gy=ghx=(gh)x\in{}O_{x}$.\n

\defn{
	We say $G$ acts \textbf{transitively} on $X$ if there is only one orbit.\n
}

\ex{
	\begin{enumerate}[leftmargin=2\parindent]
		\item $S_{n}\acts\set{1,\ldots,n}$. This action is transitive.
		\item $\GL_{n}(F)\acts{}F^{n}$. There are two orbits, $\set{0}$ and $F^{n}\cut\set{0}$.
		\item $G\acts{}G$ by left multiplication. This is transitive, because $\forall{}g\in{}G$, $g=g\cdot{}1$, so $g\in{}O_{1}$.
		\item $G\acts{}G$ by conjugation. Given $g\in{}G$, $O_{g}=\set{hgh\inv\mid{}h\in{}G}$, the set of elements conjugate to $g$.
	\end{enumerate}
}

\defn{
	The set of elements conjugate to $g$ is called the \textbf{conjugacy class} of $g$, and often denoted $C_{g}$.\n
}

\defn{
	Given $G$-sets $X$ and $Y$, we can define a new $G$-set $X\disjunion{}Y$ called the \textbf{disjoint union}, by putting together $X$ and $Y$. $X$ and $Y$ are isomorphic to $G$-subsets of $X\disjunion{}Y$ and $X\disjunion{}Y=X\cup{}Y$ and $X\cap{}Y=\emptyset$.\n
}

\prop{
	Every $G$-set is isomorphic to a disjoint union of transitive $G$-sets.\n
	Proof: $X$ is the disjoint union of its orbits.\proven
}

\defn{
	Let $X$ be a $G$-set. A subset $Y$ of $X$ is \textbf{$G$-stable} if $y\in{}Y,g\in{}G\Rightarrow{}gy\in{}Y$. $Y$ is also called a \textbf{$G$-subset}.\n
}

\defn{
	Let $X$ be a $G$-set, $x\in{}X$. The \textbf{stabilizer} of $X$ is $G_{x}=\set{g\in{}G\mid{}gx=x}$.\n
}

\par\noindent
Observe: $G_{x}$ is a subgroup of $G$.\n
Proof:
\begin{itemize}
	\item $1\in{}G_{x}$ because $1\cdot{}x=x$.
	\item If $g\in{}G_{x}$, then $gx=x$, so $g\inv{}gx=g\inv{}x$. Thus, $x=g\inv{}x$, so $g\inv\in{}G_{x}$.
	\item If $g,h\in{}G_{x}$, then $(gh)x=g(hx)=gx=x$, so $gh\in{}G_{x}$.
\end{itemize}
Thus, $G_{x}$ is a subgroup of $G$.\proven

\prop{
	Let $x\in{}X,g\in{}G$. $G_{gx}=gG_{x}g\inv$.\n
	Proof: Suppose $h\in{}G_{x}$. Then $(ghg\inv)\cdot{}gx=ghx=gx$, so $ghg\inv\in{}G_{gx}$. Thus, $G_{gx}\supseteq{}gG_{x}g\inv$.\n
	The other direction is similar.\proven
}

\cor{
	Any $2$ elements in the same orbit have conjugate stabilizers.\n
}

\par\noindent
Given $H\subseteq{}G$ a subgroup, consider the set of all left cosets $G/H=\set{gH\mid{}g\in{}G}$. This is naturally a $G$-set by $g\cdot(g'H)=gg'H$.\n

\par\noindent
Some observations:
\begin{itemize}
	\item $G/H$ is a transitive $G$-set $gH=g\cdot{}H$. Thus, every coset is in the orbit of $H$.
	\item The stabilizer of $H$ under this action is $H$, because $g\cdot{}H=gH$; if $g\in{}H$, then $g\cdot{}H=H$. Conversely, if $g\cdot{}H=H$, then $g\in{}H$.
\end{itemize}

\thm{
	Every transitive $G$-set is isomorphic to $G/H$ for some subgroup $H$. More precisely, if $x\in{}X$, $H=G_{x}$, then the function $\map{f:G/H}{X}{gH}{gx}$ is well defined and an isomorphism of $G$-sets.\n
	Proof: First, check that $f$ is well defined. Suppose $gH=g'H$. It's enough to show $gx=g'x$. We know $g=g'h$ for some $h\in{}H$. So $gx=g'hx=g'x$, because $h$ stabilizes $x$.\n
	Map of $G$-sets: Let $g,g'\in{}G$. Then $f(g\cdot(g'H))=f(gg'H0=gg'x=g(g'x)=gf(g'H)$.\n
	Surjectivity: Given $y\in{}X$, we can write $y=gx$ for some $g\in{}G_{x}$ because $X$ is transitive. Thus, $y=f(gH)$, so $y\in\im{}f$.\n
	Injectivity: Suppose $f(gH)=gx=f(g'H)=g'x$. Then $x=g\inv{}g'x$, so $g\inv{}g'\in{}G_{x}=H$, so $gH=g'H$.\proven
}

\prop{
	Given subgroups $H,H'\subseteq{}G$, $G/H$ and $G/H'$ are isomorphic as $G$-sets iff $H$ and $H'$ are conjugates.\n
	Proof: this is left as an exercise to the reader.\n
}

\cor{
	(Counting Formula) Suppose $f$ is finite, $X$ is a $G$-set, and $x\in{}X$. Then $\card{G}=\card{O_{x}}\cdot\card{G_{x}}$.\n
	Proof: $O_{x}$ is a transitive $G$-set, and $G_{x}$ is the stabilizer of $x\in{}O_{x}$. So $O_{x}\cong{}G/G_{x}$, so $\card{O_{x}}=[G:G_{x}]$.\proven
}

\end{document}